\clearpage
\VerdierPhysicalPageAnchor{0160}%
\begingroup\small\setlength{\parskip}{0pt}\setlength{\emergencystretch}{1.5em}
\VerdierSemanticAnchor{ch3-3-2-7-condition}
\noindent\textit{tel que le diagramme}
\[
X_1\xrightarrow{u_1}X_2\xrightarrow{f_3u_2}Y_3
\]
\textit{puisse s'insérer dans un triangle distingué
$(X_1,X_2,Y_3,u_1,f_3u_2,u_3)$.}
Cela résulte immédiatement de la démonstration de la proposition (3.2.5).

\VerdierSemanticAnchor{ch3-3-2-8}
\medskip\noindent\textbf{3.2.8. Démonstration des assertions a) et b) du
théorème 3.2.1.}\quad
Nous avons déjà montré (3.2.6) que la catégorie $\mathrm A(\mathcal D)$ est
abélienne. Soit $F\colon\mathcal D\to\mathcal D'$ un foncteur exact. Il
résulte immédiatement de la proposition (3.2.7) que le foncteur
$\mathrm A(F)$ est exact à gauche. Par ailleurs, il est clair d'après la
construction de la catégorie $\mathrm A(\mathcal D)$ que la catégorie
$\mathrm A(\mathcal D^\circ)^\circ$ est canoniquement isomorphe à la
catégorie $\mathrm A(\mathcal D)$ et que le foncteur
$\mathrm A(F^\circ)^\circ$ est canoniquement isomorphe au foncteur
$\mathrm A(F)$. Le foncteur $F^\circ$ étant exact (1.1.7), on en déduit que
$\mathrm A(F^\circ)$ est un foncteur exact à gauche et, par suite, que le
foncteur $\mathrm A(F)$ est aussi exact à droite, donc exact.

Les objets admissibles de $\mathcal D^\wedge$ sont isomorphes à des
sous-objets et des objets quotients de foncteurs représentables. Par suite
(3.2.3), tout objet de $\mathrm A(\mathcal D)$ est isomorphe à un sous-objet
et à un objet quotient d'un objet de $\mathcal D$. Les foncteurs
représentables sont des objets projectifs dans $\mathcal D^\wedge$. Ils sont
donc (3.2.5) projectifs dans la sous-catégorie pleine des objets admissibles
de $\mathcal D^\wedge$. Par suite (3.2.3), les objets de $\mathcal D$ sont
des projectifs de $\mathrm A(\mathcal D)$. On en déduit que les objets de
$\mathcal D^\circ$ sont des objets projectifs de
$\mathrm A(\mathcal D^\circ)$ et, par suite, en passant aux catégories
opposées, que les objets de $\mathcal D$ sont des objets injectifs de
$\mathrm A(\mathcal D)$. La démonstration en forme du fait que
$\mathrm A_{\mathcal D,\mathcal D'}$ est pleinement fidèle peut être sans
danger laissée au lecteur : un morphisme entre foncteurs exacts est
uniquement déterminé lorsqu'on connaît ses valeurs sur suffisamment d'objets
projectifs.

\VerdierSemanticAnchor{ch3-3-2-9}
\medskip\noindent\textbf{3.2.9. Fin de la démonstration du théorème 3.2.1.}
Montrons d'abord que le foncteur
$\mathcal H\colon\mathcal D\to\mathrm A(\mathcal D)$ est cohomologique. Soit
$(X,Y,Z,f,g,h)$ un triangle distingué de $\mathcal D$. Il faut montrer que la suite
\[
\mathcal H(X)\xrightarrow{\mathcal H(f)}\mathcal H(Y)
\xrightarrow{\mathcal H(g)}\mathcal H(Z).
\]
est exacte.
Par (3.2.3) et (3.2.5), il suffit de montrer que
\VerdierSemanticAnchor{ch3-3-2-9-1}
\[
h(X)\xrightarrow{h(f)}h(Y)\xrightarrow{h(g)}h(Z)
\tag{3.2.9.1}
\]
\endgroup
